\documentclass[11pt,a4paper]{article}
\usepackage[margin=27mm]{geometry}
\usepackage{amsmath,amssymb}
\usepackage[T1]{fontenc}
\setlength{\emergencystretch}{2em}
\title{Square-summable martingale noise without almost-sure convergence\\
\large Disproof of conjecture 00000008837}
\author{ziangni-sys}
\date{4 October 2026}
\begin{document}
\maketitle
The final source conjunct states that square summability of
martingale-difference noise guarantees almost-sure convergence
of stochastic splitting algorithms. Both language versions omit
solvability. A genuine backward proximal iteration on an actual
probability space disproves that unqualified implication.
The separate asserted expected convergence rate is not needed.

\textbf{Actual operator and resolvent.}
On the real Hilbert space take
\[
 A(x)=\{1\}\quad(x\in\mathbb R).
\]
The graph $G=\{(x,1):x\in\mathbb R\}$ is monotone since every
monotonicity inner product is $(x-y)(1-1)=0$.
It is maximal monotone: if a monotone graph containing $G$
contains $(u,v)$, its inequalities with $(u-1,1)$ and $(u+1,1)$
give $v-1\geq0$ and $-(v-1)\geq0$, hence $v=1$.
Its single-valued representative is globally Lipschitz.

The actual unit-parameter resolvent equation
$x\in y+A(y)$ is equivalent to $x=y+1$.
Thus every input has the unique resolvent
\[
 J_A(x)=(I+A)^{-1}(x)=x-1.
\]
Nevertheless $A$ has no zero. This missing solvability hypothesis
is explicitly part of the counterexample.

\textbf{Probability space, filtration and noise.}
Let $\Omega=\{\ast\}$ with its complete measurable space and
actual Dirac probability measure $\mu=\delta_\ast$.
For every $n\geq0$ use the trivial sigma-algebra
$\mathcal F_n=\{\varnothing,\Omega\}$.
This is an actual increasing filtration.
Set
\[
 \xi_n(\omega)=0\qquad(n\geq0,\ \omega\in\Omega).
\]
The noise is adapted and integrable, and the actual conditional
expectations satisfy
\[
 \mathbb E[\xi_{n+1}\mid\mathcal F_n]=0
\]
for every $n$. Thus it is martingale-difference noise, with no
assumed probabilistic certificate.

\newpage
\textbf{Full square summability.}
At every sample point, and in expectation, the actual squared
noise norms vanish:
\[
 \sum_{n=0}^{\infty}|\xi_n(\omega)|^2=0,\qquad
 \sum_{n=0}^{\infty}\mathbb E|\xi_n|^2=0.
\]
Both series are genuinely summable.
Consequently the counterexample satisfies either usual
pointwise or expected square-summability condition.

\textbf{Actual stochastic backward iteration.}
For an arbitrary deterministic initial point $x$, use the
standard additive-noise proximal iteration
\[
 X_0(\omega)=x,\qquad
 X_{n+1}(\omega)=J_A(X_n(\omega))+\xi_n(\omega).
\]
Every trajectory is adapted. Induction gives
\[
 X_n(\omega)=x-n
 \quad\text{for every }n\geq0,\ \omega\in\Omega.
\]
With zero noise it is also an exact backward inclusion step:
$X_n-X_{n+1}=1\in A(X_{n+1})$.
The unit proximal parameter is strictly positive; all
subproblems are defined and uniquely solved.

The actual expectation is likewise
$\mathbb E X_n=x-n$.
Allowing degenerate noise is legitimate here: the universal
square-summable-noise assertion does not require nonzero
variance or nontrivial sample space.

\textbf{No sample trajectory converges.}
If $X_n(\omega)\to a$ for a finite real $a$, the shifted sequence
has the same limit. Therefore
$X_{n+1}(\omega)-X_n(\omega)\to0$.
This difference is constantly $-1$, contradicting uniqueness
of limits.

Define the genuine convergence event
\[
 E_x=\{\omega:\exists a\in\mathbb R,\ X_n(\omega)\to a\}.
\]
The preceding argument gives $E_x=\varnothing$, hence
$\mu(E_x)=0$. Its complement has probability one.
In particular almost-sure convergence is false for every start.

\textbf{Scope and formal correspondence.}
The disproof concerns the source's unconditional stochastic
splitting assertion. It does not contradict a stochastic
proximal-point theorem that assumes a nonempty zero set.
Lean constructs the actual measure and Mathlib filtration;
checks adaptedness, integrability and actual conditional
expectations; proves both squared-noise series and actual
expectations; verifies the genuine resolvent and full recurrence;
and proves the convergence event empty and the almost-sure
statement false.

Run \texttt{lake build} in \texttt{lean/} with Lean 4.19.0 and
public Mathlib revision:
\begin{center}\small
\texttt{c44e0c8ee63ca166450922a373c7409c5d26b00b}
\end{center}
The project prints the relevant axiom audits.
\end{document}
